%%%%%%%%%%  scan idx 70  |  folio 57  %%%%%%%%%%
noté
encore $T$, de $R\times S$. On a alors $\gamma_X(R)-\gamma_X(S)=(\gamma_X(R)-\gamma_X(T))-(\gamma_X(S)-\gamma_X(T))$ et il suffit de prouver par exemple que le premier terme du second membre est dans l'image de $\underline K(Y)$. Or la projection $u:T\longrightarrow R$ est un isomorphisme sur $U$, donc $\gamma_X(R)-\gamma_X(T)=\gamma_X(\operatorname{Coker}u)-\gamma_X(\operatorname{Ker}u)$ et $\operatorname{Ker}u$ et $\operatorname{Coker}u$ ont leur support dans $Y$.

Il reste donc à prouver le lemme suivant~:

\textbf{Lemme 2.11.} \emph{Soient $X$ un espace algébrique, $F$ un faisceau algébrique cohérent sur $X$, $U$ une partie ouverte de $X$ et $G$ un sous-faisceau cohérent de $F|_U$. Alors $G$ est la restriction à $U$ d'un sous-faisceau cohérent de $F$.}

Soit $\overline G$ le sous-faisceau de $F$ dont les sections sur l'ouvert $V$ de $X$ sont les sections de $F$ sur $V$ dont la restriction à $U\cap V$ est une section de $V$. Tout revient à prouver que ce faisceau $\overline G$ est cohérent. Pour ceci, on peut supposer $X$ et $U$ affines (car $U$ est réunion d'ouverts affines $U_i$, et $\overline G$ est alors intersection des faisceaux analogues $\overline{G|_{U_i}}$). Alors l'anneau de coordonnées de $U$ est un anneau de fractions $A_S$ de l'anneau de coordonnées $A$ de $X$, et l'on vérifie aussitôt que si $F$ est défini par le $A$-module $M$, alors $F|_U$ est défini par le $A_S$-module $M_S=A_S\otimes_A M$. Le sous-faisceau $G$ de $F|_U$ est alors défini par un sous-module $N'$ de $M_S$, et l'on voit que $\overline G$ est précisément défini par le sous-module $N$ de $M$, image inverse de $N'$ par l'homomorphisme canonique de $M$ dans $M_S$.
\begin{flushright}cqfd.\end{flushright}

\textbf{Corollaire 1 à la proposition 2.10~:} \emph{Soit $x\in\underline K(X)$. Pour que $x$ soit de filtration $\leq i$, il faut et il suffit qu'il existe une partie fermée $Y$ de $X$, de dimension $\leq n-i$, telle que la restriction de $x$ à $X-Y$ soit nulle.}

